% GENERATED FILE. Edit canonical lessons, then run npm run generate:books.
% canonical-source-hash: fnv1a64:6067e6d3026e3db2
% canonical-lessons: AR-C281-raqm, AR-C281-mukalama, AR-C281-malumat, AR-C281-mudhi, AR-C281-mumaththil

\chapter{Number, Phone call, Information, Presenter, Actor}
\label{ch:ar-number-phone-call-information-presenter-actor}
\hlchaptermodality{281}

\begin{chapteropening}
\textbf{By the end of this chapter:} \emph{I can say a number, a phone call, information, a presenter and an actor.}

Complete the last lesson of chapter 281: I can say a number, a phone call, information, a presenter and an actor.
\end{chapteropening}

\section[\texorpdfstring{\ar{رقم} (raqm)}{رقم (raqm)}]{\ar{رقم} (raqm) — a number}

\label{lesson:AR-C281-raqm}



\begin{quote}
Before the new one: say the Arabic for a pilot, then the Arabic for a pharmacist.
\end{quote}

\subsection*{You'll want to know: \ar{رقم}}
\textbf{\ar{رقم}} — \emph{raqm} — \textquotedblleft{}a number\textquotedblright{}.

\begin{cousinweb}
\textbf{\ar{رقم}} (\emph{raqm}) is built on the native Arabic root \emph{r-q-m}, \textquotedblleft{}to write, to mark\textquotedblright{} (Lane). Lane's \emph{raqm} is a mark of writing, and he notes the \textquotedblleft{}Indian\textquotedblright{} numerals the Arabs adopted.
\end{cousinweb}

\begin{grammarlens}[title={What you've built}]
One more word for this chapter.
\end{grammarlens}

\subsection*{Your turn}
\begin{itemize}
\raggedright
  \item \emph{Say it:} \emph{raqm}
  \item \emph{Say it:} \emph{raqm}, once more
  \item \emph{Say it:} say \emph{raqm}
  \item \emph{Recall:} say the Arabic for a pilot, then the Arabic for a pharmacist, then say \emph{raqm} again
\end{itemize}

\begin{tcolorbox}[breakable,colback=teal!4,colframe=teal!35!black,title={Before you move on}]
What does \ar{رقم} mean? (\textquotedblleft{}A number\textquotedblright{}.) Say it once more.
\end{tcolorbox}
\section[\texorpdfstring{\ar{مكالمة} (mukālama)}{مكالمة (mukālama)}]{\ar{مكالمة} (mukālama) — a phone call}

\label{lesson:AR-C281-mukalama}



\begin{quote}
Before the new one: say the Arabic for a pharmacist, then the Arabic for a number.
\end{quote}

\subsection*{You'll want to know: \ar{مكالمة}}
\textbf{\ar{مكالمة}} — \emph{mukālama} — \textquotedblleft{}a phone call\textquotedblright{}.

\begin{cousinweb}
\textbf{\ar{مكالمة}} (\emph{mukālama}) is built on the native Arabic root \emph{k-l-m}, \textquotedblleft{}speech; to speak to\textquotedblright{} (Lane). \emph{Mukālama} is the verbal noun of Form III, \textquotedblleft{}speaking with one another\textquotedblright{}; the phone call is a modern sense.
\end{cousinweb}

\begin{grammarlens}[title={What you've built}]
One more word for this chapter.
\end{grammarlens}

\subsection*{Your turn}
\begin{itemize}
\raggedright
  \item \emph{Say it:} \emph{mukālama}
  \item \emph{Say it:} \emph{mukālama}, once more
  \item \emph{Say it:} say \emph{mukālama}
  \item \emph{Recall:} say the Arabic for a pharmacist, then the Arabic for a number, then say \emph{mukālama} again
\end{itemize}

\begin{tcolorbox}[breakable,colback=teal!4,colframe=teal!35!black,title={Before you move on}]
What does \ar{مكالمة} mean? (\textquotedblleft{}A phone call\textquotedblright{}.) Say it once more.
\end{tcolorbox}
\section[\texorpdfstring{\ar{معلومات} (maʿlūmāt)}{معلومات (maʿlūmāt)}]{\ar{معلومات} (maʿlūmāt) — information}

\label{lesson:AR-C281-malumat}



\begin{quote}
Before the new one: say the Arabic for a number, then the Arabic for a phone call.
\end{quote}

\subsection*{You'll want to know: \ar{معلومات}}
\textbf{\ar{معلومات}} — \emph{maʿlūmāt} — \textquotedblleft{}information\textquotedblright{}.

\begin{cousinweb}
\textbf{\ar{معلومات}} (\emph{maʿlūmāt}) is built on the native Arabic root \emph{ʿ-l-m}, \textquotedblleft{}to know\textquotedblright{} (Lane). \emph{Maʿlūmāt} is the plural of the passive participle, \textquotedblleft{}things known\textquotedblright{}.
\end{cousinweb}

\begin{grammarlens}[title={What you've built}]
One more word for this chapter.
\end{grammarlens}

\subsection*{Your turn}
\begin{itemize}
\raggedright
  \item \emph{Say it:} \emph{maʿlūmāt}
  \item \emph{Say it:} \emph{maʿlūmāt}, once more
  \item \emph{Say it:} say \emph{maʿlūmāt}
  \item \emph{Recall:} say the Arabic for a number, then the Arabic for a phone call, then say \emph{maʿlūmāt} again
\end{itemize}

\begin{tcolorbox}[breakable,colback=teal!4,colframe=teal!35!black,title={Before you move on}]
What does \ar{معلومات} mean? (\textquotedblleft{}Information\textquotedblright{}.) Say it once more.
\end{tcolorbox}
\section[\texorpdfstring{\ar{مذيع} (mudhīʿ)}{مذيع (mudhīʿ)}]{\ar{مذيع} (mudhīʿ) — a presenter, a broadcaster}

\label{lesson:AR-C281-mudhi}



\begin{quote}
Before the new one: say the Arabic for a phone call, then the Arabic for information.
\end{quote}

\subsection*{You'll want to know: \ar{مذيع}}
\textbf{\ar{مذيع}} — \emph{mudhīʿ} — \textquotedblleft{}a presenter, a broadcaster\textquotedblright{}.

\begin{cousinweb}
\textbf{\ar{مذيع}} (\emph{mudhīʿ}) is built on the native Arabic root \emph{dh-y-ʿ}, \textquotedblleft{}to spread abroad (news)\textquotedblright{} (Lane). \emph{Mudhīʿ} is the Form IV participle, \textquotedblleft{}one who spreads (news) abroad\textquotedblright{}; broadcasting is a modern sense.
\end{cousinweb}

\begin{grammarlens}[title={What you've built}]
One more word for this chapter.
\end{grammarlens}

\subsection*{Your turn}
\begin{itemize}
\raggedright
  \item \emph{Say it:} \emph{mudhīʿ}
  \item \emph{Say it:} \emph{mudhīʿ}, once more
  \item \emph{Say it:} say \emph{mudhīʿ}
  \item \emph{Recall:} say the Arabic for a phone call, then the Arabic for information, then say \emph{mudhīʿ} again
\end{itemize}

\begin{tcolorbox}[breakable,colback=teal!4,colframe=teal!35!black,title={Before you move on}]
What does \ar{مذيع} mean? (\textquotedblleft{}A presenter, a broadcaster\textquotedblright{}.) Say it once more.
\end{tcolorbox}
\section[\texorpdfstring{\ar{ممثل} (mumaththil)}{ممثل (mumaththil)}]{\ar{ممثل} (mumaththil) — an actor}

\label{lesson:AR-C281-mumaththil}



\begin{quote}
Before the new one: say the Arabic for information, then the Arabic for a presenter.
\end{quote}

\subsection*{You'll want to know: \ar{ممثل}}
\textbf{\ar{ممثل}} — \emph{mumaththil} — \textquotedblleft{}an actor\textquotedblright{}.

\begin{cousinweb}
\textbf{\ar{ممثل}} (\emph{mumaththil}) is built on the native Arabic root \emph{m-th-l}, \textquotedblleft{}to be like\textquotedblright{} (Lane). \emph{Mumaththil} is the Form II participle, \textquotedblleft{}one who represents\textquotedblright{}; the actor is a modern sense.
\end{cousinweb}

\begin{grammarlens}[title={What you've built}]
One more word for this chapter.
\end{grammarlens}

\subsection*{Your turn}
\begin{itemize}
\raggedright
  \item \emph{Say it:} \emph{mumaththil}
  \item \emph{Say it:} \emph{mumaththil}, once more
  \item \emph{Say it:} say \emph{mumaththil}
  \item \emph{Recall:} say the Arabic for information, then the Arabic for a presenter, then say \emph{mumaththil} again
\end{itemize}

\begin{tcolorbox}[breakable,colback=teal!4,colframe=teal!35!black,title={Before you move on}]
What does \ar{ممثل} mean? (\textquotedblleft{}An actor\textquotedblright{}.) Say it once more.
\end{tcolorbox}
